\begin{result}{B15}{Proved}{An elliptic quotient separates the genus-two quintic test}
Let $C/\mathbb F_5$ be the smooth projective curve
\[
       y^2=x^5+(x^2+1)^2,
       \qquad P=(0,1),\quad Q=\infty.
\]
Although $[P-Q]$ has exact order five in $J_C$, the points $P,Q$
give distinct geometrically prime-to-five point-section classes
over every finite extension of $\mathbb F_5$. There is a finite
\emph{\'etale} cover of degree $16$ over $\mathbb F_5$, with geometric
monodromy of order prime to five, having a rational point above $P$
and none above $Q$.
\par\emph{Passage:} cote 149, \S1, manuscript p.~17, scans 35--36,
the prime-to-characteristic covering assertion (F7).
\end{result}

\begin{proof}
Put
\[
 F(x)=x^5+(x^2+1)^2=(x-1)H(x),\qquad
 H(x)=x^4+2x^3+2x^2-x-1.
\]
Both $F$ and $H$ are square-free, and $H(1)=3$ in $\mathbb F_5$.
Thus $C$ has genus two and one rational point $Q$ at infinity.
The identity
$(y-x^2-1)(y+x^2+1)=x^5$ shows
\[
              \operatorname{div}(y-x^2-1)=5P-5Q.
\]
At $P$, the second factor is the unit $2$, while at $Q$ the
orders of $x,y$ are $-2,-5$, so that the first factor has order
$-5$. There are no other zeros or poles. The class is nonzero
because a degree-one function would make this genus-two curve
rational; hence its order is exactly five.

Take the degree-two cover $f:D\to C$ given on function fields by
\[
        v^2=x-1,\qquad w=y/v,
        \qquad w^2=H(x)=H(v^2+1).
 \tag{B15.1}
\]
The function $x-1$ has divisor $2R-2Q$, where $R=(1,0)$ is a
simple hyperelliptic branch point. The quadratic Kummer extension
is therefore unramified everywhere, including at $Q$. It is
geometrically connected: neither $x-1$ nor
$(x-1)/F(x)=1/H(x)$ is a square in $\overline{\mathbb F}_5(x)$.
The smooth projective model $D$ is hyperelliptic of genus three;
the degree-eight polynomial $H(v^2+1)$ has eight distinct roots,
since $H$ is square-free, $H(1)\ne0$, and $2$ is invertible.

The second quotient is the smooth elliptic curve
\[
                  g:D\to E,\qquad E:w^2=H(x).
\]
Choose $P'=(v,w)=(2,3)$, which maps to $P$ because
$x=v^2+1=0$ and $y=vw=1$. Its elliptic image is $r=(0,3)$.
The two points $Q_+,Q_-$ above $Q$ are rational: the polynomial
$H(v^2+1)$ has degree eight and leading coefficient one, so they
are labelled by $w/v^4=+1,-1$. The map $g$ sends them to the
two distinct rational infinity points $e_+,e_-$ of $E$, labelled
by $w/x^2=+1,-1$. In particular neither equals $r$.

Direct enumeration gives
\[
 E(\mathbb F_5)=
 \{(0,2),(0,3),(3,2),(3,3),(4,1),(4,4),e_+,e_-\}.
\]
With $r$ as origin this is an elliptic group of order eight.
Consequently each $e_\pm-r$ is nonzero and has order a power of
two. This property persists in every extension field. The norm
maps on Jacobians therefore show that neither
$[Q_\pm-P']\in J_D(\overline{\mathbb F}_5)$ is five-primary.

There is also a small actual separating cover. Let $F_E$ denote
the $5$-power Frobenius of the elliptic curve with origin $r$.
The map $F_E-1:E\to E$ is finite \emph{\'etale} of degree eight
and has constant kernel $E(\mathbb F_5)$. Form
\[
       V=D\times_{g,E,F_E-1}E\longrightarrow D\longrightarrow C,
 \tag{B15.2}
\]
where $g(P')=r$ is the identity. It has a rational point over $P'$.
It has no rational point above $Q_+$ or $Q_-$, because
$(F_E-1)(E(\mathbb F_5))=\{r\}$ whereas $g(Q_\pm)=e_\pm\ne r$.
These are all points above $Q$. The cover (B15.2) has degree
$8\cdot2=16$. Its geometric Galois closure has order prime to
five by the same deck-conjugate product argument as in B14.
Its rational fibre sets therefore distinguish the two point
sections in the specified prime-to-five quotient.

For an arbitrary finite constant extension, equality of the point
sections would, by the fixed-coset comparison in B8 and B14,
give a rational lift $Q'$ to $D$ whose section equals that of $P'$.
There are still exactly the two possibilities $Q_+,Q_-$. The
geometric prime-to-five group of $D$ is the open subgroup induced
from $C$, since the cover has degree two. Kummer theory on $J_D$
would make $[Q'-P']$ five-primary. Its nonzero two-primary norm
$e_\pm-r$ contradicts that conclusion. This proves persistence
over every finite extension, not merely separation over the
original ground field. The exact polynomial, branch and rational
point checks are recorded in \texttt{code/F7\_test5.py} and its
output; the comparison with sections uses the explicit morphisms
above.
\end{proof}
